\documentclass[11pt,a4paper]{article}
\usepackage[utf8]{inputenc}
\usepackage[T1]{fontenc}
\usepackage{amsmath}
\usepackage{amsfonts}
\usepackage{amssymb}
\usepackage{booktabs}
\usepackage{geometry}
\usepackage[numbers]{natbib}
\geometry{margin=1in}

\title{The Harmonic Foundation: Roots, Inversions, and Theoretical Context}
\author{}
\date{}

\begin{document}

\maketitle

\section{The Fundamental Root}
In modern harmonic theory, the \textbf{root} is the "parent" note of a chord. While early Renaissance theorists focused on the intervals between individual voices (the \textit{regola delle terze e seste}), a paradigm shift occurred in the 18th century.

As Jean-Philippe Rameau proposed in his seminal work, \textit{Traité de l'harmonie} (1722), chords in different inversions are functionally equivalent because they share the same root \citep{rameau1722}. This concept of the \textbf{fundamental bass} allows theorists to track the underlying harmonic progression even when the literal bass line is melodic or inverted.

\section{Triad Inversions}
An inversion occurs when a note other than the root is placed in the lowest voice. This alters the chord's stability and "flavor" without changing its functional identity.

\subsection{C Major Triad Examples $\{C, E, G\}$}
\begin{table}[h]
\centering
\begin{tabular}{@{}llll@{}}
\toprule
\textbf{Position} & \textbf{Bass Note} & \textbf{Figured Bass} & \textbf{Slash Notation} \\ \midrule
Root Position     & C (Root)           & --- (or $\binom{5}{3}$) & C                       \\
First Inversion   & E (3rd)            & $6$ (or $\binom{6}{3}$) & C/E                     \\
Second Inversion  & G (5th)            & $\binom{6}{4}$          & C/G                     \\ \bottomrule
\end{tabular}
\end{table}

\section{Seventh Chord Inversions}
With the addition of a fourth pitch class (the 7th), a third inversion becomes possible. For a G dominant seventh chord $\{G, B, D, F\}$:
\begin{itemize}
    \item \textbf{Root Position:} $G$ in bass (Symbol: $7$)
    \item \textbf{First Inversion:} $B$ in bass (Symbol: $\binom{6}{5}$)
    \item \textbf{Second Inversion:} $D$ in bass (Symbol: $\binom{4}{3}$)
    \item \textbf{Third Inversion:} $F$ in bass (Symbol: $\binom{4}{2}$ or $2$)
\end{itemize}

\section{Voice Leading and Stability}
The primary application of inversions in composition is the achievement of \textit{Fließender Gesang} (smooth melodic flow).

\begin{enumerate}
    \item \textbf{The Rule of the Shortest Way:} Inversions allow voices to move the minimum distance between chords, maintaining common tones \citep{schoenberg1911}.
    \item \textbf{Stability:} Root position triads are maximally stable. Second inversions ($\binom{6}{4}$) are treated as unstable or dissonant in classical counterpoint, requiring specific resolution, as detailed in Johann Joseph Fux's \textit{Gradus ad Parnassum} \citep{fux1725}.
\end{enumerate}

\begin{thebibliography}{9}
\bibitem{fux1725}
Fux, J. J. (1725). \textit{Gradus ad Parnassum}. (A. Mann, Trans., 1943). W. W. Norton \& Company.

\bibitem{rameau1722}
Rameau, J. P. (1722). \textit{Traité de l'harmonie réduite à ses principes naturels}. Ballard.

\bibitem{schoenberg1911}
Schoenberg, A. (1911). \textit{Theory of Harmony}. (R. E. Carter, Trans., 1978). University of California Press.

\bibitem{schoenberg1954}
Schoenberg, A. (1954). \textit{Structural Functions of Harmony}. Williams and Norgate.
\end{thebibliography}

\end{document}
